\section{The completed tower and its extension-space quotient}
\label{sec:completed-tower}

The finite markings of the preceding section supply an inverse tower
of covers of the exact-height open. Our goal is to identify its quotient
by the ordinary stabilizer with a space of independent extension
classes, including its relative structure, both marking actions, and
its compact supports. That identification is the input for the
cohomological calculations in the next part.

Fix $1\leq m<n$ and put $r=n-m$. Write $x=u_m$ and
$y_i=u_{m+i}$ for $1\leq i<r$. All analytic spaces in this
section are diamonds over $\Spa(C,\cO_C)$, with characteristic-$p$
functions obtained by tilting. We use the perfected analytic domain
\[
 X_m=\{0<|x|<1,\ |y_1|,\ldots,|y_{r-1}|<1\}
\]
associated with $A_m$; here and below an empty list of $y$-coordinates
is permitted. For a finite marking level $U\subset J$, let $Z_U\to X_m$
be the finite \'{e}tale cover obtained from $B_U$. The completed full
marking tower $Z_\infty$ is the inverse limit of these covers in
perfectoid spaces, or equivalently in diamonds. Completion always means
completion for the spectral norms on relatively compact analytic charts.
In particular, no ideal containing the invertible element $x$ is used
as an ideal of completion on $B_m$.
The chart constructions use perfected Tate algebras, rational
localizations, and fiber products of perfectoid spaces
\citep[Proposition~5.20, Theorem~6.3(i)--(ii), and
Proposition~6.18]{scholzePerfectoid2012}.

Write $\mathcal V_j=\cO(1/j)$ for the basic bundle of rank $j$ and
degree one on the Fargues--Fontaine curve. Its section sheaf is denoted
by $\mathcal H_j=H^0(\mathcal V_j)$. We use the Honda model obtained
by reducing the Lubin--Tate formal group over the unramified extension
$F_j/\Qp$ with logarithm
$\sum_{a\geq0}T^{p^{ja}}/p^a$. Its $p$-series is $T^{p^j}$.
The universal cover is a perfected open ball with Honda addition
\citep[Proposition~3.1.3]{scholzeWeinstein2013}.
The following explicit comparison fixes the endomorphism convention
as well as the underlying section sheaf.

\begin{lemma}[The covariant Honda action]
\label[lemma]{lem:tower-honda-action}
The Honda universal-cover sheaf identifies covariantly with
$\mathcal H_j$. This identification intertwines the original rational
Honda endomorphisms with the bundle endomorphisms, preserving
composition. Over a geometric base, it identifies
\[
 D_j:=\End(\mathcal V_j)\simeq\End^0(\Gamma_j),
 \qquad P_j=\cO_{D_j}^{\times}.
\]
The invariant $1/j$ uses the arithmetic-Frobenius convention in the
Honda presentation. The comparison is relative on perfectoid test
bases and is linear over the sheaf $\underline{\Qp}$ of continuous
$\Qp$-valued functions.
\end{lemma}

The relative bundle and section descriptions used below are those of
\citet[Sections~II.2--II.3]{farguesScholze2021}. The coordinate
argument proves the particular partial-height tower comparison,
including its integral frame and support data.

Put
\[
 \mathcal N_m=H^1(\mathcal V_m^{\vee}),\qquad
 (\mathcal N_m^r)_{\mathrm{ind}}
 =\{(\xi_1,\ldots,\xi_r):\xi_1,\ldots,\xi_r
       \text{ are $\Qp$-linearly independent}\}.
\]
Independence is a fiberwise condition defining the indicated open
subfunctor: relations can be normalized in the compact parameter
space $\mathbb P^{r-1}(\Qp)$, so their incidence locus has closed
image under projection. As an additive diamond, after choosing the standard untilt
description,
\begin{equation}\label{eq:tower-additive-model}
 \mathcal N_m\simeq
 \mathbb G_{a,C}^{\diamond}/\underline{F_m},
\end{equation}
where $F_m/\Qp$ is unramified of degree $m$. One obtains this by
expressing $\mathcal V_m^{\vee}$ as the pushforward of $\cO(-1)$
from the curve with coefficient field $F_m$ and applying the untilt-divisor
sequence
$0\to\cO(-1)\to\cO\to i_*C\to0$.
The resulting $H^0$ and $H^1$ sequence is
$0\to\underline{F_m}\to\mathbb G_{a,C}^{\diamond}
\to\mathcal N_m\to0$; the section and vanishing assertions are
\cite[Propositions II.2.3, II.2.5(ii), and II.3.4]{farguesScholze2021}.
This additive description does not choose affine coordinates for the
$P_m$-action: that action is the one on bundle extensions.

\begin{theorem}[Completed tower comparison]\label{thm:tower-comparison}
There is a choice of a fixed determinant-height normalization for which
\begin{equation}\label{eq:tower-quotient}
 [Z_\infty/G]\simeq(\mathcal N_m^r)_{\mathrm{ind}}.
\end{equation}
The map $Z_\infty\to(\mathcal N_m^r)_{\mathrm{ind}}$ is a
$G=P_n$-torsor. The comparison commutes with base change and with both
integral marking actions of
$J=P_m\times\GL_r(\Zp)$. Under the convention that a matrix acts
on a kernel framing by pushout, $(a,b)\in P_m\times\GL_r(\Zp)$
acts on extensions by pullback along $a^{-1}$ on $\mathcal V_m$ and
pushout along $b$ on $\cO^r$.
\end{theorem}

The proof separates three geometric constructions.  The normalized
coordinates first identify the completed tower with only the \'{e}tale
marking as $(\mathcal H_n^r)_{\mathrm{ind}}$
(\cref{lem:tower-etale-comparison}).  Independent sections then give
an exact sequence with quotient of type $\mathcal V_m$, and the
connected marking supplies an actual relative frame of that quotient.
Finally the determinant height selects the integral frame torsors.
We first establish the Honda-action comparison needed to retain the
original groups throughout these constructions.

\begin{proof}[Proof of \cref{lem:tower-honda-action}]
Work on an affinoid perfectoid test base $T=\Spa(R,R^+)$ over $\kk$.
The chosen inclusion $\mathbb F_{p^j}\subset\kk$
induces the unramified field $F_j$ in the relative period coefficients;
Witt Frobenius $\phi$ restricts to arithmetic Frobenius there.
For a topologically nilpotent universal-cover coordinate $X$, the
covariant Lubin--Tate section map is
\begin{equation}\label{eq:tower-honda-period-map}
 \psi_j(X)=\sum_{i\in\Z}p^i[X^{p^{-ji}}].
\end{equation}
Here brackets denote Teichm\"uller lifts. This is the natural additive
isomorphism of \citet[Proposition~II.2.2]{farguesScholze2021}, applied
with coefficient field $F_j$, uniformizer $p$, and residue-field
cardinality $p^j$. Pushforward along the unramified degree-$j$ map
of coefficient-field curves gives
$\cO(1/j)=f_*\cO(1)$, as in the proof of
\citet[Theorem~II.2.14]{farguesScholze2021}. Thus its target is
$\mathcal H_j$, and $\phi^j\psi_j(X)=p\psi_j(X)$.

We record the continuity needed to use this comparison in families.
On a closed period annulus, its Gauss seminorms can be normalized by
$|p|=p^{-1}$ and $|[z]|_s=|z|^s$, with $s$ in a compact interval
$[a,b]$ and $a>0$. On an input chart $\|X\|\leq\rho<1$, the
$i$th term of \eqref{eq:tower-honda-period-map} has norm at most
\[
 p^{-i}\rho^{a p^{-ji}}.
\]
For $i\to+\infty$ the first factor forces convergence. For $i=-h$
the bound is $p^h\rho^{a p^{jh}}$, which also tends to zero.
Both tails converge uniformly on the input chart. The formula is
therefore continuous, commutes with bounded maps of perfectoid
algebras, and glues on test bases. Termwise calculation gives
\[
 \psi_j(X^p)=\phi\psi_j(X),\qquad
 \psi_j(cX)=[c]\psi_j(X)\quad(c\in\mathbb F_{p^j}).
\]
Consequently the Honda Frobenius $\Pi_j:X\mapsto X^p$ and the
Teichm\"uller scalars satisfy, on both sides with the same composition,
\[
 \Pi_j^j=p,\qquad \Pi_j[c]=[c^p]\Pi_j.
\]
These are the Honda algebra relations of
\citet[A2.2.17--A2.2.18]{ravenel1986}. They identify the original
integral endomorphism ring as a finite free $\Zp$-algebra generated
by $\Pi_j$ and a primitive unramified Teichm\"uller root. In
particular finitely many monomials in these generators span it over
$\Zp$. Compatibility with integers extends to $\Zp$ by continuity:
the discrepancy after an approximation modulo $p^h$ has Honda norm
at most $\rho^{p^{jh}}$, and its period-side norm tends to zero on
each annulus. The displayed action identities therefore hold for
the entire original maximal order, and then for its rational division
algebra after inverting $p$.

For the required sheaf linearity, let
$\alpha\in C^0(|T|,\Qp)$ and fix an input section $v$.
Quasicompactness supplies $d\geq0$ such that $b=p^d\alpha$ is
$\Zp$-valued. Choose its finite clopen-constant approximation $b_h$
modulo $p^h$, and write $b-b_h=p^hc_h$, with $c_h$ integral.
Integral Honda scalar operations do not increase the norm, so
\[
 \|[b]v-_{\Gamma_j}[b_h]v\|\leq\rho^{p^{jh}}.
\]
On each fixed period annulus, multiplication by an integral scalar
has norm at most one, and
\[
 \|(b-b_h)\psi_j(v)\|\leq |p|^h\|\psi_j(v)\|.
\]
On each clopen piece constant scalar compatibility gives
$\psi_j([b_h]v)=b_h\psi_j(v)$. Taking the two limits and then
commuting with $p^{-d}$ proves the assertion for $\alpha$ on the
same test base.

Positive-slope vanishing identifies derived sections of
$\mathcal V_j$ with its ordinary section sheaf. Relative full
faithfulness \citep[Corollary~3.11]{anschutzLeBras2025}, with the
just-verified topological-field sheaf convention, turns the Honda
action into a unital, composition-preserving map
$\End^0(\Gamma_j)\to\End(\mathcal V_j)$. Its source is a division
algebra, so the map is injective. Over a geometric base both algebras
have dimension $j^2$
over $\Qp$: for the target, the endomorphism bundle has slope zero
and rank $j^2$, and geometric classification together with
$H^0(\cO)=\Qp$ computes its sections. Thus the map is an
isomorphism. The unique maximal orders correspond, identifying the
actual Honda automorphism group with the bundle-frame group.
Uniqueness in full faithfulness makes this construction compatible
with relative base change. On a general test base the corresponding
action is by the sheaf $\underline{D_j}$ of continuous $D_j$-valued
functions; its global sections need not have finite dimension over
$\Qp$. No opposite action is introduced.
\end{proof}

\subsection{Normalized division coordinates}

Set $q_n=p^n$ and $q_m=p^m$.  We begin with the universal cover of a
fixed deformation over a bounded perfectoid chart.  Here the parameters
$u_i$ are given; later they will themselves be functions of the tuple
of division coordinates.  In a coordinate reducing to the Honda
coordinate, the multiplication series of the universal formal group
$\mathcal F$ satisfies
\begin{equation}\label{eq:tower-series-bounds}
 [p]_{\mathcal F}(T)\equiv T^{q_n}
       \pmod{(u_m,\ldots,u_{n-1})},\qquad
 [p]_{\mathcal F}(T)\in A_m[[T^{q_m}]].
\end{equation}
All coefficients are integral. Thus on a chart on which
$\max |u_i|\leq\lambda<1$, and for inputs of norm less than one,
uniformly on each closed chart of strictly subunit input radius,
\[
 |[p]_{\mathcal F}(s)-s^{q_n}|\leq\lambda,
 \qquad
 |[p]_{\mathcal F}(s)-[p]_{\mathcal F}(t)|
       \leq |s-t|^{q_m}.
\]
The second inequality is a coefficientwise power-series estimate in
the coordinate $T^{q_m}$, so it also holds for uniform norms over a
relative perfectoid affinoid. Every fixed-level series evaluation
converges on its own chart of radius less than one. The coefficientwise
Lipschitz constant is one, independently of that radius, so these
estimates remain uniform when the radii of successive finite
approximants tend to one; the parameter bound $\lambda<1$ stays fixed.

For a compatible division sequence
$[p]_{\mathcal F}t_{j+1}=t_j$, define
\begin{equation}\label{eq:tower-normalized-limit}
 z=\lim_{j\to\infty}t_j^{q_n^j}.
\end{equation}
Consecutive terms differ by at most $\lambda^{q_n^j}$; hence this
limit converges uniformly. Conversely, for $|z|<1$, the formula
\begin{equation}\label{eq:tower-inverse-cover}
 t_j=\lim_{N\to\infty}
       [p]_{\mathcal F}^{N}
          \bigl(z^{1/q_n^{j+N}}\bigr)
\end{equation}
converges uniformly: consecutive approximants differ by at most
$\lambda^{q_m^N}$. The formulas are inverse. Indeed, the difference
between $t_l$ and $z^{1/q_n^l}$ has norm at most $\lambda$, and
applying $[p]^{l-j}$ bounds the difference at level $j$ by
$\lambda^{q_m^{l-j}}$, which tends to zero. This proves both
uniqueness and recovery of a given compatible sequence.  Conversely,
for the sequence constructed from a given $z$, finite telescoping
gives $|t_j-z^{1/q_n^j}|\leq\lambda$, and hence
\[
 |t_j^{q_n^j}-z|\leq\lambda^{q_n^j}\longrightarrow0.
\]
Thus the other composite is the identity as well.  It also shows
that the kernel of projection to $t_0$ is the integral Tate module of
the \'{e}tale part: on a perfect geometric fiber, the connected torsion
has no nonzero points. Projection to $t_0$ is v-locally surjective,
since the monic division equations admit roots after a finite cover
at every level.

We now let the deformation vary with the division tuple.  The whole
reduced closures of \cref{prop:mark-roots} give free polydisc
coordinates at each finite level; this extra input will identify the
completed tower itself.

\begin{lemma}[The analytic \'{e}tale-marking tower]
\label[lemma]{lem:tower-etale-comparison}
Let $Z_\infty^{\mathrm{et}}$ denote the completed perfected tower of
ordered \'{e}tale Tate-module markings over $X_m$.  There is a relative
analytic isomorphism
\[
 Z_\infty^{\mathrm{et}}\simeq(\mathcal H_n^r)_{\mathrm{ind}},
\]
where independence is fiberwise over $\Qp$.  The isomorphism is
equivariant for $G$ and $\GL_r(\Zp)$.  It is the restriction of an
analytic isomorphism from the completed perfected tower of reduced
closures $R_{m,l}$ to the open ball underlying $\mathcal H_n^r$.
\end{lemma}

\begin{proof}
Adjoining the point coordinates to the powered \'{e}tale coordinates
is a pure-root extension by \cref{prop:mark-roots}; it does not change
a perfected analytic level.  Thus these point coordinates compute
the completed \'{e}tale marking tower.  Use first the reduced
level-$l$ closure in \cref{prop:mark-roots}. Its coordinate ring has
regular parameters
$t_{1,l},\ldots,t_{r,l}$. Two consequences of that proposition are
\begin{align}
 u_i&\in k[[t_{1,1}^{q_m},\ldots,t_{r,1}^{q_m}]],
       \quad u_i(0)=0,\label{eq:tower-first-order}\\
 t_{i,j}&=F_{i,j,l}
      (t_{1,l}^{q_m^{l-j}},\ldots,t_{r,l}^{q_m^{l-j}})
       \quad(j\leq l),\label{eq:tower-early-coordinates}
\end{align}
with $F_{i,j,l}$ integral convergent formal series, without a constant
term. Formula \eqref{eq:tower-early-coordinates} follows also by
iterating the second assertion of \eqref{eq:tower-series-bounds}
and expressing the base parameters in the powered level-$l$
coordinates.

For $l\geq1$ and $0<\rho<1$ in the divisible value group, take the closed chart
\[
 X_l(\rho)=
 \{\max_i |t_{i,l}^{q_n^l}|\leq\rho\}
\]
in the perfected generic fiber of this regular ring. If
$\lambda=\max |u_i|$ at a point of this chart, then
\eqref{eq:tower-first-order} and finite telescoping of
\eqref{eq:tower-series-bounds} give
\[
 \lambda\leq\|t_{\cdot,1}\|^{q_m},\qquad
 \|t_{\cdot,1}\|
    \leq\max(\lambda,\rho^{1/q_n}).
\]
If $\lambda>\rho^{1/q_n}$ these inequalities would give
$\lambda\leq\lambda^{q_m}<\lambda$. Therefore
\begin{equation}\label{eq:tower-uniform-parameter}
 \lambda\leq\rho^{q_m/q_n}<1.
\end{equation}
For $y_{i,l}=t_{i,l}^{q_n^l}$ this yields
\begin{equation}\label{eq:tower-normalized-error}
 \|y_{\cdot,l+1}-y_{\cdot,l}\|
   \leq\lambda^{q_n^l}
   \leq\rho^{q_mq_n^{l-1}}<\rho.
\end{equation}
The finite integral injective transition maps are surjective on
these charts. To check the radius condition in this assertion, lift a
point by lying over in the finite map. The bound
\eqref{eq:tower-uniform-parameter} holds at the original point and
therefore at its lift; \eqref{eq:tower-normalized-error} then forces
the lifted normalized coordinate to have norm at most $\rho$.
Consequently pullback on functions is isometric for spectral norms,
also after adjoining all $p$-power roots.

Let $L_\rho$ be the completed direct limit of these perfected
affinoid algebras. The limits $z_i=\lim_l y_{i,l}$ define a map from
the perfected radius-$\rho$ Gauss algebra to $L_\rho$. This map is
isometric. To see this without inferring it from points, take a
polynomial $f$ in finitely many $p$-power roots of the variables. At
every sufficiently late level, $f(y_{1,l},\ldots,y_{r,l})$ has
exactly the radius-$\rho$ Gauss norm of $f$: the $t_{i,l}$ are free
regular polydisc coordinates, and taking powers followed by perfection
does not change the perfected Gauss algebra. These evaluations converge
to $f(z_1,\ldots,z_r)$ in the common spectral norm. The limit has the
same norm, since the differences eventually have smaller norm than a
given nonzero Gauss norm. Completion gives the asserted isometry; its
image is therefore closed.

It is surjective as well. For fixed $i,j$, replace $t_{a,l}$ in
\eqref{eq:tower-early-coordinates} by $z_a^{1/q_n^l}$. The
substituted series converges on the radius-$\rho$ polydisc and the
error is bounded by
\begin{equation}\label{eq:tower-surjectivity-error}
 \lambda^{q_m^{l-j}}
 \leq\rho^{(q_m/q_n)q_m^{l-j}}\longrightarrow0.
\end{equation}
Thus $t_{i,j}$ is in the closed image. Taking any fixed $p$-power
root in this argument merely takes that root of the error bound, so
all perfected finite-level generators are in the image. They are
dense in $L_\rho$, proving surjectivity. This establishes inverse
affinoid morphisms. Varying $\rho$ gives the isomorphism of open
balls.

All the preceding estimates are uniform norm estimates for integral
series. Hence the constructed morphisms commute with bounded base
change, glue on overlapping charts, and apply with any compact
profinite parameter space $\mathcal D$ by using the uniform norm on
$C(\mathcal D,C^\flat)$. In particular they establish a relative
analytic comparison, not just a comparison of geometric points.

The normalized limits respect addition, integral scalar operations,
and the full $G$-action. For example, replace the deformation addition
series by its Honda reduction. Their difference on integral inputs
has norm at most $\lambda$. Transporting to a fixed earlier division
level contracts this error by $q_m^{l-j}$, exactly as in
\eqref{eq:tower-surjectivity-error}. The same argument applies to
the integral coordinate change associated with $g\in G$ and its
Honda reduction; compact families of $g$ have the same coefficient
bounds. It applies also to the finite sums defining integral changes
of an \'{e}tale basis. Thus the comparison is equivariant for $G$
and $\GL_r(\Zp)$, continuously in both groups.

We can now identify the exact-height open. An integral relation among
the normalized sections gives, by the inverse construction, the same
relation among their compatible division sequences: if the normalized
relation is zero, its point at a fixed earlier level is bounded by
$\lambda^{q_m^{l-j}}$ for every $l$, and is zero. The converse follows
by applying the limit formula. Clearing denominators gives the same
statement for rational relations. Over height exactly $m$, the marked
\'{e}tale Tate module has rank $r$ and its marked basis is independent.
Over a deeper height its rank is less than $r$, so the compatible
tuple is dependent. It follows that restricting this analytic isomorphism to the
exact-height open gives
\begin{equation}\label{eq:tower-etale-independent}
 Z_\infty^{\mathrm{et}}\simeq(\mathcal H_n^r)_{\mathrm{ind}}.
\end{equation}
Here the ambient analytic morphisms have already been constructed;
the geometric-point criterion is used only to identify their open
subfunctors. The original basis condition is open and closed at each
finite level above $x\ne0$, so no additional basis component is
lost in this restriction.
\end{proof}

\subsection{Subbundles, extensions, and relative frames}

\begin{proof}[Completion of the proof of \cref{thm:tower-comparison}]
We translate \eqref{eq:tower-etale-independent} into bundle
geometry. The geometric classification and the relative constant
Harder--Narasimhan theorem used in this step are
\cite[Theorems II.2.14 and II.2.19]{farguesScholze2021}.

Let $s_1,\ldots,s_r$ be independent sections of $\mathcal V_n$
over a geometric point, and let $\mathcal W$ be the saturation of
their generic span. It is generically generated by $\cO^r$, so it
has no negative-slope quotient: a map from $\cO$ to a negative-slope
bundle is zero. If $b=\rk\mathcal W\leq r<n$, stability of
$\mathcal V_n$ gives
$\deg\mathcal W<b/n<1$. Its degree is a nonnegative integer,
so it is zero. Slope classification now gives
$\mathcal W\simeq\cO^b$. The map from $\cO^r$ to this bundle is
a matrix over $\Qp$. Independence forces $b=r$ and the matrix to
be invertible onto its image. This proves the subbundle assertion
over every geometric base curve.

For a general perfectoid test base $T$, write $X_T$ for its relative
Fargues--Fontaine curve and $Y_T$ for the covering space with Frobenius
quotient $X_T$. Let $i_T:\cO^r\to\mathcal V_n$ be the actual map
supplied by the sections.
Its dual $i_T^\vee:\mathcal V_n^\vee\to\cO^r$ is surjective on every
geometric base curve. To see that these curves detect every point,
take $x\in X_T$ and lift it to $y\in Y_T$, whose Frobenius quotient
is $X_T$.
Choose a geometric point $\Spa(L^\sharp,L^{\sharp+})$ over $y$.
The identification $Y_T^\diamond=T\times\operatorname{Spd}(\Qp)$
makes its tilt a map $\Spa(L,L^+)\to T$. The untilt $L^\sharp$ gives
a point of $Y_L$ mapping to $y$, hence a point of the geometric base
curve $X_L$ mapping to $x$
\cite[Definition II.1.15 and Propositions II.1.16--II.1.18]{farguesScholze2021}.
Thus at every $x$ some maximal minor of $i_T^\vee$ is nonzero in
the residue field: otherwise it would remain zero at the point of
$X_L$, contradicting surjectivity there. That minor is a unit in the
stalk and on a neighborhood, so $i_T^\vee$ is surjective and splits
locally. Dualizing makes $i_T$ a locally split injection with locally
free cokernel $\mathcal Q$. Its formation commutes with base change.
We have therefore obtained the relative exact sequence
\begin{equation}\label{eq:tower-bundle-sequence}
 0\longrightarrow\cO^r\longrightarrow\mathcal V_n
   \longrightarrow\mathcal Q\longrightarrow0.
\end{equation}
On each geometric fiber, every slope of $\mathcal Q$ is positive:
a nonpositive-slope quotient
would give a forbidden map from the positive stable bundle
$\mathcal V_n$. It has degree one and rank $m$. Each positive basic
summand has positive integral degree, so there is exactly one such
summand, of degree one; hence $\mathcal Q\simeq\mathcal V_m$.
Thus the already constructed vector bundle $\mathcal Q$ has constant
basic type. The relative classification
\cite[Theorem II.2.19(ii), including the dual-surjection argument
in its proof]{farguesScholze2021}
now supplies its local frames.

Conversely, an extension
\begin{equation}\label{eq:tower-extension}
 0\longrightarrow\cO^r\longrightarrow\mathcal E
    \longrightarrow\mathcal V_m\longrightarrow0
\end{equation}
has no negative-slope quotient, by applying the same vanishing of
maps separately to the two ends. A nonnegative bundle of degree one
is positive basic unless it has a quotient $\cO$. The connecting map
\[
 \Hom(\cO^r,\cO)=\Qp^r
    \longrightarrow\Ext^1(\mathcal V_m,\cO)=\mathcal N_m
\]
sends $(a_i)$ to $\sum_i a_i\xi_i$, where $(\xi_i)$ is the
extension class. Since $\Hom(\mathcal V_m,\cO)=0$, a quotient
$\mathcal E\to\cO$ exists exactly when this map has a nonzero
kernel. Thus the middle term is $\mathcal V_n$ exactly on
$(\mathcal N_m^r)_{\mathrm{ind}}$. Moreover
$\Hom(\mathcal V_m,\cO^r)=0$, so an extension with both ends fixed
has no nonidentity automorphism inducing the identity on the ends.
This identifies the extension moduli with a sheaf, with no residual
automorphism group.

It remains to show that an actual connected formal marking supplies
the claimed quotient frame in families. Its recursively monic equations
in \cref{prop:mark-jets} bound all its coefficients by one on the
charts above. Its series therefore converges on the open ball and
induces a map on universal covers
$\mathcal H_n\to\mathcal H_m$. This is a morphism of sheaves of
$\underline{\Qp}$-modules, where
$\underline{\Qp}(T)=C^0(|T|,\Qp)$. Indeed it commutes with addition
and the integral scalar operations, and multiplication by $p$ is
invertible on universal covers. It therefore commutes with every
constant rational scalar. Here is the continuity argument for a
varying scalar on the same affinoid perfectoid test object $T$.
Write this section map as $\psi$, fix $v\in\mathcal H_n(T)$, and
let $\alpha\in C^0(|T|,\Qp)$. Compactness gives an integer $a\geq0$
such that $b=p^a\alpha$ takes values in $\Zp$. Let $b_k$ be the
finite-valued, clopen-constant representative of $b$ modulo $p^k$,
and put $c_k=(b-b_k)/p^k$. In the fixed Honda coordinates choose
$\|v\|\leq\rho<1$ and $\|\psi(v)\|\leq\rho'<1$. Integral scalar
operations do not increase these norms. Since
$[p^k]_{\Gamma_j}(T)=T^{p^{jk}}$, the Honda differences satisfy
\[
 \|[b]v-_{\Gamma_n}[b_k]v\|
   =\|[p^k c_k]v\|\leq\rho^{q_n^k},\qquad
 \|[b]\psi(v)-_{\Gamma_m}[b_k]\psi(v)\|
   \leq(\rho')^{q_m^k}.
\]
Both bounds tend to zero. On each clopen piece,
$\psi([b_k]v)=[b_k]\psi(v)$ by constant scalar compatibility.
The continuity of $\psi$, established by the convergent analytic
formulas, therefore gives $\psi([b]v)=[b]\psi(v)$. Commuting with
$p^{-a}$ proves the same identity for $\alpha$ on $T$. This proves
linearity for the topological-field sheaf without changing the test
object.

The marking kills the compatible integral division sections, since
their images are Honda torsion on a perfect characteristic-$p$ base,
and hence kills their $\Qp$-span. The relative section sequence of
\eqref{eq:tower-bundle-sequence} is exact, since
$H^1(\cO)=0$ as a v-sheaf. Thus
$\mathcal H_n\to H^0(\mathcal Q)$ is an epimorphism of
$\underline{\Qp}$-module sheaves. The map factors through that
epimorphism, and its factor remains $\underline{\Qp}$-linear.

The relative section functor on perfect complexes is fully faithful
by \cite[Corollary 3.11, pp. 3680--3681]{anschutzLeBras2025},
applied over the present small v-stack. The preceding construction
gives a morphism in $D(T_v,\underline{\Qp})$, compatibly on all test
bases, with the topological-field convention of that source.
Positive bundles have no
higher relative cohomology sheaves, so this statement applies to
the ordinary section sheaves just constructed. It gives a unique
bundle map
\[
 f:\mathcal Q\longrightarrow\mathcal V_m.
\]
This map is nonzero on every geometric fiber. Indeed, the linear
coefficient of the formal marking is nonzero on $x\ne0$, and on a
sufficiently small ball it dominates the higher terms. Universal
cover projection is locally surjective, so the induced section map
is nonzero there. Both bundles on a geometric fiber have the same
stable basic type, whose endomorphisms form a division algebra.
Thus $f$ is an isomorphism on every fiber, and hence is a relative
bundle isomorphism by the same pointwise maximal-minor criterion.
Full faithfulness and uniqueness make this
construction compatible with every base change and with the three
group actions. This is the needed relative frame construction.

\subsection{The integral height of a frame}

For a frame of $\mathcal Q$, use the ordered exact sequence
\eqref{eq:tower-bundle-sequence} to identify
$\det\mathcal Q$ with $\det\mathcal V_n$, and fix the standard
identifications
$\det\mathcal V_n\simeq\cO(1)\simeq\det\mathcal V_m$.
The determinant of the frame is then a section of the locally
profinite sheaf $\underline{\Qp^{\times}}$. Its valuation is a
locally constant integer. The reduced-norm valuation
\[
 v_p\circ\operatorname{Nrd}:D_m^{\times}\longrightarrow\Z
\]
is surjective and has kernel $P_m$. Therefore frames of any specified
height form a $P_m$-torsor, locally nonempty.

The height of the frame obtained from a connected marking is unchanged
by $P_m$. An integral \'{e}tale basis change multiplies the ordered
determinant by a unit, so it also preserves this height. The height
therefore descends from the full marking tower through both marking
torsors to $X_m$. The domain $X_m$ is connected: it is exhausted by
connected annulus-times-polydisc charts with nonempty successive
intersections; perfection preserves their underlying spaces. Hence
the descended height is one constant integer $c$.

Changing the fixed determinant normalization absorbs $c$. Over
$(\mathcal H_n^r)_{\mathrm{ind}}$, the connected-marking tower is
a $P_m$-torsor, and so is the space of quotient frames of this
normalized height. The frame map is equivariant for that same group.
An equivariant map between two torsors for the same group is an
isomorphism: after a local section of the source, its image is a
section of the target, and the map becomes the identity on the
group coordinate. This proves surjectivity onto every frame of the
specified height as well as injectivity.

Finally start with an arbitrary independent extension
\eqref{eq:tower-extension}. Its middle term has type $\mathcal V_n$,
so its rational frames form a $D_n^{\times}$-torsor. The determinant
of the framed ends specifies a height for a middle frame. Surjectivity
of the reduced-norm valuation makes the locus of that height locally
nonempty, and it is a $P_n$-torsor. A frame in this locus produces
\eqref{eq:tower-bundle-sequence} together with a quotient frame of
the normalized height. By the preceding paragraph this is precisely
a point of $Z_\infty$. Thus every independent extension occurs and
its middle-frame fiber is exactly $G=P_n$. Taking the quotient gives
\eqref{eq:tower-quotient}. Integral changes at either end, and the
ordinary stabilizer on the middle frame, change determinant only by
a unit. They preserve the selected heights and act by the stated
pushout and pullback operations. More explicitly, for an extension
with maps $(i,q)$, postcomposing the connected marking by $a$ changes
$q$ to $aq$; the resulting class is the pullback along $a^{-1}$.
Under the stated kernel-frame convention, $b$ changes $i$ to
$ib^{-1}$ and gives pushout by $b$. These use the covariant Honda
action of \cref{lem:tower-honda-action} and commute with the
middle-frame action. This completes the proof of
\cref{thm:tower-comparison}.
\end{proof}

\subsection{Transport of supports}

For later use we specify the support convention. A compact-support
complex is the filtered colimit of restriction fibers away from
relatively compact quasi-compact bounds. Closed boxes can be
interleaved with slightly larger open boxes. On $X_m$ one can take
\[
 \epsilon\leq|x|\leq\rho<1,\qquad |y_i|\leq\rho,
\]
with separated endpoint radii and with $\epsilon\downarrow0$,
$\rho\uparrow1$. Supports over a compact profinite parameter space
$\mathcal D$ must fit one such bound uniformly in $\mathcal D$.
The same convention is used on finite marking covers and on their
completed towers.

\begin{lemma}[Compact torsors and support bounds]
\label[lemma]{lem:tower-supports}
The two compact torsors
\[
 Z_U\ \longleftarrow\ Z_\infty
       \longrightarrow\ (\mathcal N_m^r)_{\mathrm{ind}},
 \qquad\text{with groups }U\text{ and }G,
\]
transport the preceding systems of support bounds cofinally. For
the characteristic-$p$ function sheaf, and also with uniform compact
profinite parameters, descent identifies
\begin{equation}\label{eq:tower-support-descent}
 \RGamma\bigl(G,\RGamma_c(Z_U,\cO^\flat)\bigr)
 \simeq
 \RGamma\bigl(U,\RGamma_c(
       (\mathcal N_m^r)_{\mathrm{ind}},\cO^\flat)\bigr).
\end{equation}
The group cochain terms retain the support convention uniformly in
their compact bar parameters.
\end{lemma}

\begin{proof}
After pullback to a perfectoid test space, a torsor for a locally
profinite group is represented by a pro-\'{e}tale, universally open
v-cover \cite[Lemma 10.13]{scholzeDiamonds2017}. For a profinite group
$H$, that lemma presents it as $\varprojlim_K P/K$, where $K$ ranges
over open subgroups and $P/K$ is finite \'{e}tale over the test space.
Over an affinoid perfectoid test space these finite covers are
affinoid, and their limit is affinoid perfectoid, hence quasi-compact.
Quasi-compactness descends from these test covers to the torsor
morphism by \cite[Proposition 10.11(o)]{scholzeDiamonds2017}.
Thus the same assertion holds for the locally spatial diamonds here;
in particular a quasi-compact bound has quasi-compact inverse image.
Conversely a quasi-compact bound upstairs has a quasi-compact
image. The chosen increasing, slightly open box exhaustion covers
the space; any such image has a finite subcover by these boxes and
is contained in one larger box. The same argument after a finite
refinement of charts works on finite covers. The actual relative
isomorphism of \cref{thm:tower-comparison} transports all
quasi-compact bounds, so these observations prove the asserted
cofinality on both sides. Compact parameter spaces do not change
the argument: include the parameter in the quasi-compact bound, or
use a finite clopen refinement and its uniform norm.

On an invariant bound, the \v{C}ech nerve of a compact-group torsor
is its continuous bar construction. Supports pull back to the
specified inverse-image bound in every degree. Every bound upstairs
can be enlarged to a bound invariant under both groups: its saturation
by the compact product $G\times U$ is quasi-compact and is contained
in a larger bound. Descent for the function sheaf and for the restriction fiber
therefore gives the continuous group-cochain description on these
cofinal invariant bounds. The two actions on $Z_\infty$ commute,
so taking first $U$-descent and then $G$-descent, or in the opposite
order, gives the same double bar complex. The support and sheaf
complexes are bounded below. In each total degree only finitely
many bar and sheaf degrees contribute, and filtered colimits commute
with these finite sums and restriction fibers. We may consequently
pass to the support colimit in the double bar complex, obtaining
\eqref{eq:tower-support-descent}. This also proves its assertion
about all compact bar parameters. The noncompact additive cover in
\eqref{eq:tower-additive-model} is not used in this compact-torsor
argument; supports for that cover are computed separately in
\cref{sec:kummer,sec:independent-tuples}.
\end{proof}
